\subsection{Certification of the running-cost arithmetic}
\label{sec:cost-roundoff}

This subsection concerns the frozen functions \texttt{B}, \texttt{Bnode}, and
\texttt{log\_lower}; it does not introduce a different running-cost bound.
We use the constants $\varepsilon_{\mathrm{in}}=2^{-48}$,
$\varepsilon_{\mathrm{ar}}=2^{-32}$, and $\mathfrak u=2^{-53}$ fixed above.
Write $\mathfrak g_j=j\mathfrak u/(1-j\mathfrak u)$.
Hats denote computed binary64 quantities.  In particular,
$|\widehat p-p|<8\varepsilon_{\mathrm{in}}$ and
$|\widehat t-t|<2\varepsilon_{\mathrm{in}}$, as established in the
continuation analysis.  The validated data satisfy
\[
 \begin{gathered}
 \Delta=1/640,\quad 0\le t\le1-\Delta,\quad 0\le x<y\le1,\\
 \tfrac14<b<2,\quad |p|<2,\quad 0\le I_r\le\Delta,\quad 0\le E(r),e^{-r}\le1.
 \end{gathered}
\]
Here $p=cx+y(D-b)$, where $.85<c<.9$ and $0\le D\le1$: thus
$p>-2$ and $p\le y(.9+1-1/4)\le1.65<2$.
Each stored coefficient has absolute error less than
$\varepsilon_{\mathrm{in}}$, and the stored $\Delta$ has relative error
less than $\varepsilon_{\mathrm{in}}$.  Write
\[
 p_-=-(-p)^+,\qquad p_+=p^+,\qquad
 C(r)=p_-\Delta+p_+I_r,
\]
\[
 \alpha(r)=\left(ybE(r)-\frac{p_-I_r}\Delta\right)\Delta^2,
 \quad \beta(r)=b(1-y+ty)\Delta+ybI_r\Delta,
 \quad z(r)=\beta(r)e^{-r}.
\]

\paragraph{A reproducible error bound for the coefficients.}
For every elementary operation, including one with a subnormal result, we
use the following weaker form of the preceding floating-point model:
\begin{equation}
 |\operatorname{fl}(v)-v|\le\mathfrak u|v|+2^{-1074}.
 \label{eq:cost-basic-roundoff}
\end{equation}
For example, if $|a|\le A$, $|b|\le B$ and the operand errors are
$e_a,e_b$, a product has error at most
\[
 Ae_b+Be_a+e_ae_b+
 \mathfrak u(A+e_a)(B+e_b)+2^{-1074}.
\]
An addition has error at most
$e_a+e_b+\mathfrak u(A+B+e_a+e_b)+2^{-1074}$.
For division by a positive exact $b\ge b_0>e_b$, first bound the input
perturbation by $(e_a+|a/b|e_b)/(b_0-e_b)$, and then apply
\eqref{eq:cost-basic-roundoff}.  The maps $v\mapsto\min(v,0)$ and
$v\mapsto\max(v,0)$ do not increase absolute errors.

The following table applies these inequalities in the source's evaluation
order.  The middle column bounds the magnitude of the corresponding exact
expression, and the last column is a strict error bound in units of
$\varepsilon_{\mathrm{in}}$.  Each multiplication by $\Delta$ is a
separate operation.  The dyadic quantity $1-y$ is exact.
\begin{center}
\small
\begin{tabular}{lll}
\hline
Expression & Exact magnitude bound & Error coefficient\\
\hline
$yb$; $ybE$ & $2$; $2$ & $2$; $5$\\
$p_-I_r$; $p_-I_r/\Delta$ & $2\Delta$; $2$ & $3$; $2000$\\
$ybE-p_-I_r/\Delta$ & $4$ & $2010$\\
$(ybE-p_-I_r/\Delta)\Delta$ & $4\Delta$ & $4$\\
$\alpha=(ybE-p_-I_r/\Delta)\Delta^2$ & $4\Delta^2$ & $1$\\
$ty$; $1-y+ty$ & $1$; $1$ & $3$; $4$\\
$b(1-y+ty)$; $b(1-y+ty)\Delta$ & $2$; $2\Delta$ & $10$; $1$\\
$ybI_r$; $ybI_r\Delta$ & $2\Delta$; $2\Delta^2$ & $3$; $1$\\
$\beta$; $z=\beta e^{-r}$ & $2\Delta(1+\Delta)<.004$ & $3$; $4$\\
$p_-\Delta$; $p_+I_r$; $C=p_-\Delta+p_+I_r$ & $2\Delta$ & $1$; $3$; $5$\\
\hline
\end{tabular}
\end{center}
For the largest intermediate error, for example, division of the preceding
$3\varepsilon_{\mathrm{in}}$ error by
$\Delta(1-\varepsilon_{\mathrm{in}})$ contributes less than
$1921\varepsilon_{\mathrm{in}}$; denominator perturbation and rounding
contribute less than $3\varepsilon_{\mathrm{in}}$.  The two following
multiplications reduce the bound to the displayed $4$ and $1$.
All other entries follow by the same displayed product and sum rules;
the strict margins also exceed the sum of the fewer than $100$ possible
$2^{-1074}$ terms.  In particular, all four final errors are less than
$8\varepsilon_{\mathrm{in}}$.

\paragraph{The logarithm computed by the frozen source.}
This paragraph certifies the actual twelve-term loop, not the unused
degree-five polynomial mentioned in an older explanation.  Whenever the
logarithmic branch is reached, put
$v=\widehat\beta/\widehat\alpha$ in exact real arithmetic and
$\widehat v=\operatorname{fl}(v)$.  The guards give
$\widehat\beta>\widehat\alpha>\varepsilon_{\mathrm{ar}}$; hence
$1\le\widehat v<2^{25}$ and $|\log\widehat v-\log v|<2\mathfrak u$.
The exact binary operation \texttt{frexp} gives $\widehat v=2^e m$ with
$0\le e\le24$ and $1\le m<2$.  For $w_0=(m-1)/(m+1)$,
\[
 S_{12}(m)=2\sum_{j=0}^{11}\frac{w_0^{2j+1}}{2j+1}
 \le\log m,\qquad 0\le w_0<\tfrac13.
\]
The subtraction $m-1$ is exact.  Rounding of $m+1$ and the division give
a relative error bounded by $\mathfrak g_2$
in a nonzero $w_0$.  In the largest monomial, its degree is $23$;
the repeated squaring-factor and product operations contribute at most
$22$ further factors, division contributes one, and summation contributes
at most eleven.  Thus the computed $2\,\texttt{sum}$ is at most
\[
 (1+\mathfrak g_{128})S_{12}(m)+256\cdot2^{-1074}.
\]
This bound also covers underflow of the smallest positive monomials.
There is no assumption about a platform logarithm implementation.

The hexadecimal constant in the source is the exact rational
\[
 c_2=\frac{3121657384082679}{2^{52}}
       =\texttt{0x1.62e42fefa39eep-1}.
\]
The finite rational inequality
\[
 2\sum_{j=0}^{39}\frac1{(2j+1)3^{2j+1}}
       > c_2+2^{-53}
\]
is checked by multiplication by the positive common denominator.
The left side is a lower partial sum of the positive series for $\log2$,
so $c_2<\log2$.  Since $S_{12}(m)<1$ and $e\le24$, the above
polynomial error, the product $ec_2$, the final addition and subtraction,
and the rounded quotient have combined upper error at most
\[
 \mathfrak g_{128}S_{12}(m)+24\mathfrak u+25\mathfrak u
 +26\mathfrak u+2\mathfrak u+512\cdot2^{-1074}
 <256\mathfrak u+512\cdot2^{-1074}
 <32\varepsilon_{\mathrm{in}}.
\]  Consequently the returned value satisfies
\begin{equation}
 \texttt{log\_lower}(\widehat v)
 \le\log v-\varepsilon_{\mathrm{ar}}
            +32\varepsilon_{\mathrm{in}}
 <\log v.
 \label{eq:cost-log-certified}
\end{equation}

\paragraph{The free bound and its branch tests.}
Recall $\mathcal H$ from~\eqref{eq:cert-iid-free}.
On the guarded domain and on the segment joining the exact and computed
coefficient pairs, $\alpha>\varepsilon_{\mathrm{ar}}/2$ and
$\beta<.004$.  The function is continuously differentiable across
$\beta=\alpha$, and direct differentiation gives
\[
 0\le\partial_\beta \mathcal H\le1,\qquad
 0\le\partial_\alpha \mathcal H<10.
\]
Indeed $\beta/\alpha<2^{26}$ in the logarithmic region.  Thus coefficient
perturbations change $\mathcal H$ by less than
$88\varepsilon_{\mathrm{in}}$.  This argument permits the computed and
exact pairs to fall on different sides of the branch test.
In the rational branch the division, subtraction, and multiplication
contribute less than $\varepsilon_{\mathrm{in}}$ in absolute error.
In the other branch, \eqref{eq:cost-log-certified} and the remaining two
rounded operations give the same one-sided bound.
Including the error of \texttt{base}, addition, and subtraction, the
computed candidate \texttt{base+freebound-eps} is therefore at most
\[
 C(r)+\mathcal H(\alpha(r),\beta(r))
       -\varepsilon_{\mathrm{ar}}+128\varepsilon_{\mathrm{in}},
\]
and is a valid lower bound.  If a positivity guard fails, the source simply
does not use this candidate.  The baseline \texttt{min(p,0)*dt-eps}
is separately a lower bound for $\min(p,0)\Delta$.

\paragraph{The scalar shifted bound.}
The frozen rate grid contains $64$ exactly and every other rate has
distance at least $16$ from $64$.  Thus \texttt{rr[hi]<=64} cannot admit
an exact rate exceeding $64$.  The test
$\widehat z>\operatorname{fl}(\widehat\alpha+
\varepsilon_{\mathrm{ar}})$ implies $z>\alpha$: the total coefficient
and addition errors are less than $17\varepsilon_{\mathrm{in}}$.
The same strict inequality holds throughout the segment between the
exact and computed $(\alpha,z)$ pairs.  For
\[
 \mathcal J(\alpha,z,r)=z+\alpha r/2-
       \frac{\alpha^2}{24(z-\alpha/2)},\qquad z>\alpha\ge0,
\]
differentiation gives
\[
 1\le\partial_z\mathcal J\le\tfrac76,
 \qquad |\partial_\alpha\mathcal J|\le33,
 \qquad |\partial_r\mathcal J|=\alpha/2<1
 \quad(0\le r\le64+\varepsilon_{\mathrm{in}}).
\]
Hence coefficient perturbations, including that of $r$, contribute less
than $275\varepsilon_{\mathrm{in}}$.  The subtraction in
$\ell=z-\alpha/2$ has no harmful relative cancellation, since
$\ell>z/2$ and $\ell>\alpha/2$.
The exact quotient $\alpha^2/(24\ell)$ is at most $\alpha/12$.
The rounding of $\ell$, its multiplication by $24$, the numerator
product, and division give absolute error at most
$\mathfrak g_5\alpha/12+2^{-1000}$ in this quotient.
For the additive term, the guard gives $\ell>\varepsilon_{\mathrm{ar}}/2$;
propagating the fewer than ten possible $2^{-1074}$ errors through these
operations amplifies them by less than $2^{40}$, which is covered by
$2^{-1000}$.  All remaining products and sums have magnitude less
than $1$; applying \eqref{eq:cost-basic-roundoff} to them and to
\texttt{base} bounds the complete upward error by
$512\varepsilon_{\mathrm{in}}$.
Thus \texttt{sharp-eps} is below the exact shifted expression whenever
this branch is used.

\paragraph{The shifted bound for a tree node.}
Let the node's endpoint rates be $r_-\le r_+\le64$.
The quantities actually formed by the source are
\[
 \overline\alpha=\operatorname{fl}(\widehat\alpha(r_-)
                 +16\varepsilon_{\mathrm{in}}),\quad
 \underline\alpha=\max\{0,\operatorname{fl}(\widehat\alpha(r_+)
                 -16\varepsilon_{\mathrm{in}})\},
\]
\[
 \underline z=\max\{0,\operatorname{fl}(\widehat z(r_+)
                 -16\varepsilon_{\mathrm{in}})\},\quad
 \underline r=\max\{0,\operatorname{fl}(\widehat r_-
                 -16\varepsilon_{\mathrm{in}})\}.
\]
The table above proves
$\overline\alpha\ge\alpha(r_-)$,
$0\le\underline\alpha\le\alpha(r_+)$, and
$0\le\underline z\le z(r_+)$: the extra additions and subtractions
have error less than $\varepsilon_{\mathrm{in}}$.
For the last expression their error is at most
$3\varepsilon_{\mathrm{in}}$, so $0\le\underline r\le r_-$ as well.
The rounded guard implies $\underline z>\overline\alpha$.
Therefore the exact expression in these already directed floating inputs,
\[
 C(r_+)+\underline z+
 \underline\alpha\,\underline r/2-
 \frac{\overline\alpha^2}
      {24(\underline z-\overline\alpha/2)},
\]
is a lower bound for the shifted cost of every descendant original row,
by the endpoint argument proved above.
Its quotient is again at most $\overline\alpha/12$, with absolute
operation error at most $\mathfrak g_5\overline\alpha/12+2^{-1000}$.  The remaining operation errors and
the error of \texttt{base} together are less than
$16\varepsilon_{\mathrm{in}}$.  Subtracting
$2\varepsilon_{\mathrm{ar}}$ therefore preserves the lower-bound
direction.  No monotonicity of the full shifted expression is assumed.

In \texttt{B}, the maximum of the available certified candidates remains
certified.  The final subtraction of $\varepsilon_{\mathrm{ar}}$, exact
power-of-two scaling, floor, and subtraction of one integer unit give a
lower bound in units of $2^{-40}$.  The analogous floor and subtraction
in \texttt{Bnode} have the same direction.  The magnitudes just established
are below $1$, so the final floating subtraction has error less than
$\varepsilon_{\mathrm{in}}$ and cannot exhaust any of these margins.

\subsection{Table magnitudes, integer safety, and the final objective}
\label{sec:integer-safety}

The allowance computation requires no approximation theorem.  The
coordinates are exact multiples of $2^{-30}$, so its $w=y-x$ is exact
and lies in $[0,1]$.  Whatever error occurs in computing the intermediate
\texttt{future}, the two comparisons in
\texttt{max(0,min(w,w-future))} put the returned floating value in $[0,w]$.
Multiplication by $2^{24}$ is exact, and consequently
\[
 0\le\texttt{allowance}\le2^{24}.
\]
Subtracting this integer therefore only lowers a proposed table value,
by at most $2^{-16}$.  The fixed additional slack is $2^{-28}$.

Let $M_k$ be the largest absolute table entry in layer $k$, after division
by $2^{40}$.  All entries are nonpositive.  Integer barycentric
interpolation of the next layer lies between its minimum and zero:
flooring cannot cross the integer minimum.  The previously certified
multiplier, including its ceiling and extra integer unit, is at most
$2^{40}(1+2\varepsilon_{\mathrm{ar}}+2\cdot2^{-40})$.
The subsequent integer floor loses less than $2^{-40}$.
For the running cost, the actual baseline arithmetic gives
\[
 2^{-40}\texttt{B(hi)}
 \ge -3/640-3\varepsilon_{\mathrm{ar}}-2\cdot2^{-40}.
\]
Indeed $|\widehat p|<2+8\varepsilon_{\mathrm{in}}$;
the computed product with $\widehat\Delta$ is greater than $-3/640$;
the baseline and final floating subtractions use two
$\varepsilon_{\mathrm{ar}}$ deductions, their rounding is covered by
the third, and flooring followed by subtraction of one costs less than
two integer units.  Taking a maximum can only improve this lower bound.
It follows, with deliberately enlarged constants, that
\begin{align*}
 M_k&\le aM_{k+1}+d,\qquad M_{640}=0,\\
 a&=1+2\varepsilon_{\mathrm{ar}}+2\cdot2^{-40},\\
 d&=3/640+4\varepsilon_{\mathrm{ar}}+8\cdot2^{-40}
                  +2^{-28}+2^{-16}.
\end{align*}
Taking minima with zero or with additional proposals, and either
monotonicity-lowering pass, preserves this global lower bound.
Since $640(a-1)<1$, the binomial theorem gives
\[
 M_k\le d\sum_{j=0}^{639}a^j
 \le\frac{640d}{1-640(a-1)}
 =\frac{1846693964}{613566573}<\frac{301}{100}<4.
\]
Thus every stored integer has magnitude below $4\cdot2^{40}$.

Whenever interpolation is performed, its tests give
$0\le X<Y\le2^{30}$.  Lookup indices lie in $[0,65536]$, and every
grid width and offset lies in $[0,2^{30}]$.  The barycentric denominator
and each unscaled weight are at most $2^{60}$, so their products and
differences fit signed 64-bit integers.  The selected triangle has
nonnegative weights summing to the denominator.  Its weighted numerator
and all partial sums therefore have magnitude at most
$4\cdot2^{40}\cdot2^{60}=2^{102}$, safely within signed 128-bit range.
Continuation multiplier products are smaller than $2^{84}$, and the
negative-floor formula's additional denominator is also harmless at
these bounds.  Each conversion of a running-cost value, multiplier, or
allowance to a signed 64-bit integer is within its range.
The state-array size is $1201\cdot1202/2=721801$; all index products fit
signed 32-bit integers.  A binary tree over $131093$ leaves has fewer
than twice that many nodes; the additional proposal evaluations are
fewer than $50$ per state.  Per-state counters fit unsigned 32-bit
integers and their sums over all states fit unsigned 64-bit integers.

Finally, input validation establishes
$\texttt{HU}=\lfloor2^{40}\theta(0)\rfloor$.
The reported exact objective is therefore
\[
 \frac{\texttt{HU}+\texttt{table}[m,0]}{2^{40}}
 \le\theta(0)+\frac{\texttt{table}[m,0]}{2^{40}}.
\]
Together with the original-row feasibility induction, this proves the
claimed direction of the rational objective; the printed decimal plays
no role in acceptance.
